\article{轨道积分与基变换}{Orbital integrals and base change}{R. Kottwitz}{2}{111--113}

设 $F$ 是特征为 $0$ 的局部域（local field），$E$ 或者为 $F\times\cdots\times F$（$l$ 个因子），或者为 $F$ 的 $l$ 次循环扩张（cyclic extension），其中 $l$ 是素数。将 $F$ 对角嵌入 $F\times\cdots\times F$。记 $\mathcal O_F$ 为 $F$ 的赋值环（valuation ring）；当 $E=F\times\cdots\times F$ 时，令 $\mathcal O_E=\mathcal O_F\times\cdots\times\mathcal O_F$；当 $E$ 是域时，令 $\mathcal O_E$ 为 $E$ 的赋值环。若 $E$ 是域，令 $\Gamma=\Gal(E/F)$，并取 $\Gamma$ 的一个生成元 $\sigma$。若 $E=F\times\cdots\times F$，令 $\sigma$ 为 $E$ 的自同构
\[
 (x_1,\ldots,x_l)\longmapsto(x_2,\ldots,x_l,x_1),
\]
并令 $\Gamma$ 为由 $\sigma$ 生成的 $E$ 的自同构群；此时 $\Gamma$ 同样是 $l$ 阶循环群。

令 $G=\GL_2$。$\Gamma$ 在 $E$ 上的作用诱导 $\Gamma$ 在 $G(E)$ 上的作用。借助 $F$ 到 $E$ 的嵌入，将 $G(F)$ 与 $G(E)^\Gamma$ 视为同一群。定义范数映射（norm map）$N:G(E)\to G(E)$ 为
\[
 Ng=g^{\sigma^{l-1}}\cdots g^\sigma g.
\]
这一映射由 Saito [2] 引入，它依赖于 $\sigma$ 的选择。

\begin{originaltheorem}{引理 1（Lemma 1）}
设 $g,x\in G(E)$，则
\begin{align*}
\text{(i)}\quad &N(g^{-\sigma}xg)=g^{-1}(Nx)g;\\
\text{(ii)}\quad &(Nx)^\sigma=x(Nx)x^{-1};\\
\text{(iii)}\quad &\det(Nx)=N_{E/F}(\det x);\\
\text{(iv)}\quad &Nx\text{ 在 }G(E)\text{ 中共轭于 }G(F)\text{ 的某个元素。}
\end{align*}
\end{originaltheorem}

前三个断言可直接计算验证，(iv) 也可由 (ii) 推出。受等式 (i) 启发，作如下定义：若存在 $g\in G(E)$ 使 $y=g^{-\sigma}xg$，便称 $G(E)$ 中的 $x,y$ 为 $\sigma$-共轭（$\sigma$-conjugate）。

(i) 与 (iv) 合起来表明，$N$ 诱导一个从 $G(E)$ 的 $\sigma$-共轭类到 $G(F)$ 的共轭类的映射。这一映射是单射；若 $E=F\times\cdots\times F$，它还是满射。还应注意，$x$ 与 $y$ 为 $\sigma$-共轭，当且仅当 $(\sigma,x)$ 与 $(\sigma,y)$ 在半直积（semidirect product）$\Gamma\ltimes G(E)$ 中共轭。

分别在 $G(F)$ 和 $G(E)$ 上选择 Haar 测度（Haar measure）$dg,dg_E$。若 $F$ 是非 Archimedes 的（nonarchimedean），则将测度正规化，使
\[
 \operatorname{meas}(K_F)=\operatorname{meas}(K_E)=1,
 \qquad K_F=G(\mathcal O_F),\quad K_E=G(\mathcal O_E).
\]
对 $G(F)$ 的任意元素 $\gamma$，用 $G_\gamma$ 表示 $\gamma$ 在 $G$ 中的中心化子（centralizer）。下面给出 Shintani [4] 提出的定义。\footnote{原文的 AMS（MOS）主题分类（1970）：主要分类 22E50。原文版权：1979，American Mathematical Society。}

\begin{originaltheorem}{定义（Definition）}
设 $\varphi\in C_c^\infty(G(E))$，$f\in C_c^\infty(G(F))$。若满足以下条件，便称 $f$ 与 $\varphi$ 相伴（associated）：当 $N\delta=\gamma$ 且 $\gamma$ 为 $G(F)$ 的正则元素（regular element）时，
\begin{equation}
 \int_{G_\gamma(F)\backslash G(F)}f(g^{-1}\gamma g)\frac{dg}{dt}
 =\int_{G_\gamma(F)\backslash G(E)}\varphi(g^{-\sigma}\delta g)\frac{dg_E}{dt};
 \tag{A}
\end{equation}
对于 $G(F)$ 中不能写成 $G(E)$ 中元素的范数的正则元素 $\gamma$，则要求
\begin{equation}
 \int_{G_\gamma(F)\backslash G(F)}f(g^{-1}\gamma g)\frac{dg}{dt}=0.
 \tag{B}
\end{equation}
在 (A)、(B) 中，$dt$ 是 $G_\gamma(F)$ 上的 Haar 测度。
\end{originaltheorem}

\paragraph*{注记（Remarks）.}
(1) 对于每个是范数的正则元素 $\gamma\in G(F)$，只需对一个 $\delta$ 检查定义中的条件 (A)。

(2) 对任意 $\delta\in G(E)$，令
\[
 G_\delta^\sigma(E)=\{g\in G(E):g^{-\sigma}\delta g=\delta\}.
\]
不难看出，若 $\gamma=N\delta$ 是 $G(F)$ 的正则元素，则 $G_\delta^\sigma(E)=G_\gamma(F)$。对 $\gamma\in G(F)$，若 $\gamma$ 是中心元素，而且它是 $G(E)$ 中某个不与中心元素 $\sigma$-共轭的元素的范数，则定义 $\varepsilon(\gamma)=-1$；其他情形定义 $\varepsilon(\gamma)=1$。可以证明，若 $f$ 与 $\varphi$ 相伴，则当 $N\delta=\gamma\in G(F)$ 时，
\begin{equation}
 \int_{G_\gamma(F)\backslash G(F)}f(g^{-1}\gamma g)\frac{dg}{dt}
 =\varepsilon(\gamma)\int_{G_\delta^\sigma(E)\backslash G(E)}
 \varphi(g^{-\sigma}\delta g)\frac{dg_E}{dt'};
 \tag{A$'$}
\end{equation}
这里 $dt'$ 是 $G_\delta^\sigma(E)$ 上的 Haar 测度，它由 $G_\gamma(F)$ 上的 Haar 测度 $dt$ 确定。此外，若 $\gamma\in G(F)$ 不是 $G(E)$ 中元素的范数，则有
\begin{equation}
 \int_{G_\gamma(F)\backslash G(F)}f(g^{-1}\gamma g)\frac{dg}{dt}=0.
 \tag{B$'$}
\end{equation}

\begin{originaltheorem}{引理 2（Lemma 2）}
(i) 给定 $\varphi\in C_c^\infty(G(E))$，至少存在一个与 $\varphi$ 相伴的 $f\in C_c^\infty(G(F))$。

(ii) 给定 $f\in C_c^\infty(G(F))$，$f$ 能与某个 $\varphi\in C_c^\infty(G(E))$ 相伴，当且仅当对 $G(F)$ 中每个不是来自 $G(E)$ 的范数的元素 $\gamma$ 都有
\[
 \int_{G_\gamma(F)\backslash G(F)}f(g^{-1}\gamma g)\frac{dg}{dt}=0
\]
（等价地，只要求对每个不是来自 $G(E)$ 的范数的正则元素 $\gamma\in G(F)$ 成立）。
\end{originaltheorem}

假设 $F$ 是非 Archimedes 的。记 $\mathcal H_F=\mathcal H(G(F),K_F)$ 为 $G(F)$ 上复值、紧支集且关于 $K_F$ 双不变的函数所成的 Hecke 代数（Hecke algebra）。令 $\mathcal H_E=\mathcal H(G(E),K_E)$。Saito [2] 引入了一个 $\C$-代数同态（$\C$-algebra homomorphism）$b:\mathcal H_E\to\mathcal H_F$，下面说明其定义。

对 $f\in\mathcal H_F$，令 $f^\vee(\pi)=\Tr\pi(f)$（$\pi\in D_F$），这样便在 $G(F)$ 的非分歧不可约容许表示（unramified irreducible admissible representation）的同构类集合 $D_F$ 上定义了函数 $f^\vee$。集合 $D_F$ 是 $\C$ 上的代数簇（algebraic variety）：它同构于 $\C^*\times\C^*$ 在对称群（symmetric group）$S_2$ 交换两个因子的作用下的商。映射 $f\mapsto f^\vee$ 给出 $\mathcal H_F$ 到 $D_F$ 上正则函数代数的一个 $\C$-代数同构，称为 Satake 同构（Satake isomorphism）。上述讨论同样适用于 $E$：若 $E$ 是域，则 $D_E$ 仍然同构于 $(\C^*\times\C^*)/S_2$；若 $E=F\times\cdots\times F$，则 $D_E=D_F\times\cdots\times D_F$。

存在一个从 $D_F$ 到 $D_E$ 的代数簇映射：若 $E=F\times\cdots\times F$，它是 $\pi\mapsto\pi\otimes\cdots\otimes\pi$；若 $E$ 是域，它是
\[
 \pi(\tau)\longmapsto\pi\bigl(\Res^{W_F}_{W_E}\tau\bigr),
\]
其中 $\pi(\tau)$ 是 $G(F)$ 中对应于 $F$ 的 Weil 群（Weil group）$W_F$ 的非分歧表示 $\tau$ 的非分歧表示，而 $\pi(\Res^{W_F}_{W_E}\tau)$ 是 $G(E)$ 中对应于 $\tau$ 在 $E$ 的 Weil 群 $W_E$ 上的限制的非分歧表示。沿这一映射拉回正则函数，就得到从 $D_E$ 的正则函数代数到 $D_F$ 的正则函数代数的 $\C$-代数同态，再经 Satake 同构得到 $b:\mathcal H_E\to\mathcal H_F$。

\paragraph*{注记.}
若 $E=F\times\cdots\times F$，则
\[
 \mathcal H_E\simeq\mathcal H_F\otimes\cdots\otimes\mathcal H_F,
 \qquad b(f_1\otimes\cdots\otimes f_l)=f_1*\cdots*f_l.
\]

\begin{originaltheorem}{引理 3（Lemma 3）}
若 $F$ 是非 Archimedes 的，且 $E$ 在 $F$ 上非分歧（unramified）（允许 $E=F\times\cdots\times F$），则对所有 $\varphi\in\mathcal H_E$，$b(\varphi)$ 都与 $\varphi$ 相伴。
\end{originaltheorem}

当 $E=F\times\cdots\times F$ 时，这个引理容易证明。当 $E$ 是域时，Saito [2] 首先证明了它；随后 Langlands [1] 又利用 $\SL_2(F)$ 与 $\SL_2(E)$ 的建筑（building）给出了证明。

令 $A$ 为 $\GL_2$ 中的对角矩阵群。整体域（global field）上 $\GL_2$ 的迹公式（trace formula）包含与 $A$ 的正则元素对应的加权轨道积分（weighted orbital integral）。需要在 $G(F)$ 上引入一个函数 $\lambda_F$，其对数就是这些积分中的权因子。设 $g\in G(F)$，将其写成
\[
 g=a\begin{pmatrix}1&x\\0&1\end{pmatrix}k,
 \qquad a\in A(F),\quad k\in K_F.
\]
当 $x\in\mathcal O_F$ 时令 $\lambda_F(g)=1$，否则令 $\lambda_F(g)=|x|^{-2}$。当 $E$ 为域时，以同样方式在 $G(E)$ 上定义函数 $\lambda_E$。

\begin{originaltheorem}{引理 4（Lemma 4）}
假设 $F$ 是非 Archimedes 的，且 $E$ 是 $F$ 的非分歧扩域。则对任意满足 $N\delta$ 正则的 $\delta\in A(E)$ 以及任意 $\varphi\in\mathcal H_E$，有
\[
 l\int_{A(F)\backslash G(F)}b(\varphi)\bigl(g^{-1}(N\delta)g\bigr)
 \log\lambda_F(g)\frac{dg}{da}
 =\int_{A(F)\backslash G(E)}\varphi(g^{-\sigma}\delta g)
 \log\lambda_E(g)\frac{dg_E}{da},
\]
其中 $da$ 是 $A(F)$ 上的任意 Haar 测度。
\end{originaltheorem}

证明见 [1] 的第 3 节。

\begin{thebibliography}{4}
\bibitem[1]{c26-1} R. P. Langlands, \emph{Base change for GL(2)}, Notes, Institute for Advanced Study, 1975.
\bibitem[2]{c26-2} H. Saito, \emph{Automorphic forms and algebraic extensions of number fields}, Lectures in Mathematics, Kyoto University, 1975.
\bibitem[3]{c26-3} H. Saito, \emph{Automorphic forms and algebraic extensions of number fields}, U.S.-Japan Seminar on Number Theory (Ann Arbor, 1975).
\bibitem[4]{c26-4} T. Shintani, \emph{On liftings of holomorphic automorphic forms}, U.S.-Japan Seminar on Number Theory (Ann Arbor, 1975)（亦见本论文集第二册，第 97--110 页）。
\end{thebibliography}

\begin{flushleft}\small 美国华盛顿州 Richland\end{flushleft}
